% !PRINTED 103
\weberpage{103}
\begin{center}
\textbf{§ 31. \gesperrt{The linear fundamental transformations
of the $\vartheta$-functions.}}
\end{center}

In the application to the transformation of the $\Theta$-functions
[§ 27, (12)] the transformed modulus
\[
\omega' = \dfrac{c + \partial\omega}{a + b\omega}
\]
first comes into consideration. This does not alter if the four transformation numbers
receive a common factor $m$. The
transformation is called a \gesperrt{proper} one if $a, b, c, \partial$ have no
common factor.

The transformed argument
\[
u' = \dfrac{nu}{a + b\omega}
\]
passes over into $mu'$ if $a, b, c, \partial$ are replaced by $ma, mb, mc, m\partial$, and hence $n$
by $m^2n$. The multiplication $\begin{pmatrix}m, & 0\\0, & m\end{pmatrix}$ leaves the
modulus $\omega$ unchanged and changes $u$ into $mu$.

According to the results of the two preceding sections, the whole system of proper
transformations can be derived by repeated application of the three transformations
\[
\begin{pmatrix}1, & 0\\1, & 1\end{pmatrix}, \quad \begin{pmatrix}0, & 1\\-1, & 0\end{pmatrix}, \quad \begin{pmatrix}n, & 0\\0, & 1\end{pmatrix},
\]
and so we consider first the linear fundamental transformations
of the $\vartheta$-functions.

I. $\begin{pmatrix}1, & 0\\1, & 1\end{pmatrix}$ or $(\omega, \omega + 1)$.

By § 27, (11), (12) we have
\weq{1}{\vartheta_{11}(u, \omega + 1) = A\vartheta_{11}(u, \omega),}
in which $A$ is independent of $u$. If one replaces $u$ by
\[
u + \tfrac{1}{2}, \quad u + \tfrac{\omega}{2}, \quad u + \tfrac{1+\omega}{2},
\]
then the formulae § 21, (8) give
\weq{2}{\vartheta_{10}(u, \omega + 1) = A\vartheta_{10}(u, \omega)}
\weq{3}{e^{\frac{\pi i}{4}}\vartheta_{00}(u, \omega + 1) = A\vartheta_{01}(u, \omega)}
\weq{4}{e^{\frac{\pi i}{4}}\vartheta_{01}(u, \omega + 1) = A\vartheta_{00}(u, \omega).}
